% from ega0-21-fr.tex
\oldpage[0\textsubscript{IV-1}]{164}
Soit $K'$ le sous-corps (égal au sous-anneau) de $K$ engendré par
$k(K^p)$ et les $x_\alpha$~; compte tenu de (21.4.3), il suffit de voir que,
pour chaque $\alpha$, $x_\alpha$ n'appartient pas à $K_\alpha$. Or, dans la
suite exacte
\[
\Omega^1_{K_\alpha/k(K^p)}\otimes_{K_\alpha}K\longrightarrow
\Omega^1_{K/k(K^p)}\longrightarrow\Omega^1_{K/K_\alpha}\longrightarrow0,
\]
les images par la flèche de gauche des $d_{K_\alpha/k}(x_\beta)$ pour
$\beta\neq\alpha$ sont les $d_{K/k}(x_\beta)$, donc engendrent un sous-espace
vectoriel de $\Omega^1_{K/k(K^p)}$ ne contenant pas $d_{K/k}(x_\alpha)$~; comme
ce sous-espace est le noyau de la flèche de droite, on voit que
$d_{K/K_\alpha}(x_\alpha)\neq0$, donc $x_\alpha\notin K_\alpha$.

Si maintenant $(x_\alpha)$ est une famille $p$-libre, elle est une sous-famille
d'une $p$-base de $K$ sur $k$ (21.4.2)~; les $d_{K/k}(x_\alpha)$ font donc
partie d'une base de l'espace vectoriel $\Omega^1_{K/k}$ (21.2.5), et sont
par suite linéairement indépendants sur $K$. Réciproquement, si les
$d_{K/k}(x_\alpha)$ sont linéairement indépendants sur $K$, les deux
paragraphes précédents montrent que $(x_\alpha)$ est $p$-libre sur $k$.
\end{proof}

\begin{corollary}[21.4.6]
\label{0.21.4.6-fr}
Soient $k$ un corps de caractéristique $p>0$, $K$ une extension de $k$,
$x$ un élément de $K$. Les trois conditions suivantes sont équivalentes~:
\begin{enumerate}
  \item[a)] $x\notin k(K^p)$.
  \item[b)] $d_{K/k}(x)\neq0$.
  \item[c)] L'élément $x$ est $p$-libre sur $k$.
\end{enumerate}
\end{corollary}

\begin{proposition}[21.4.7]
\label{0.21.4.7-fr}
Soient $L$ un corps de caractéristique $p>0$, $K$ un sous-corps de $L$ tel
que $L^p\subset K\subset L$. Alors la suite d'espaces vectoriels sur $L$
\begin{equation}
\tag{21.4.7.1}\label{0.21.4.7.1-fr}
0\to\Omega^1_{K/L^p}\otimes_KL\to\Omega^1_{L/L^p}
\to\Omega^1_{L/K}\to0
\end{equation}
est exacte~; en d'autres termes, on a $\Upsilon_{L/K/L^p}=0$.
\end{proposition}
\begin{proof}
C'est un cas particulier de (21.2.4), compte tenu de (21.4.2).
\end{proof}

\begin{proposition}[21.4.8]
\label{0.21.4.8-fr}
Sous les hypothèses de (21.4.7), l'homomorphisme canonique
\begin{equation}
\tag{21.4.8.1}\label{0.21.4.8.1-fr}
(\Omega^1_{L/K})^{(p)}\longrightarrow\Upsilon_{L/K}
\end{equation}
est bijectif~; si $(x_\lambda)$ est une $p$-base de $L$ sur $K$, les éléments
$d_K(x_\lambda^p)\otimes1$ forment une base de l'espace vectoriel $\Upsilon_{L/K}$.
\end{proposition}
\begin{proof}
C'est un cas particulier de (21.3.5).
\end{proof}

\subsection{Application~: critères de séparabilité}
\label{0.21.5-fr}

Dans ce numéro et les deux suivants, on ne suppose plus que les anneaux
considérés soient de caractéristique $p>0$.

\begin{env}[21.5.1]
\label{0.21.5.1-fr}
Notons d'abord que le critère (21.2.7) permet de démontrer une partie du
théorème de Cohen sur les extensions séparables (19.6.1), savoir que si $k$
est de caractéristique $p>0$ et si $K$ est une extension séparable de $k$,
alors $K$ est une $k$-algèbre formellement lisse. En effet, $K$ admet une
$p$-base sur $k$ (21.4.2), et d'autre part, il résulte du critère de MacLane
(Bourbaki, \emph{Alg.}, chap.~V, §~8, n$^\circ$~2, prop.~3) que, dans une
clôture algébrique de $K$, $k^{p^{-1}}$ et $K$ sont linéairement disjoints sur
$k$, et par suite que l'homomorphisme canonique
$k^{p^{-1}}\otimes_k K\to k^{p^{-1}}(K)$ est bijectif, ce qui est précisément
la condition (i) de (21.2.7), après transport de structure par
l'isomorphisme $x\mapsto x^p$.
\end{env}
\begin{proposition}[21.5.2]
\label{0.21.5.2-fr}
\textnormal{(MacLane).}
Soient $A$, $B$ deux anneaux de valuation discrète complets, $\mathfrak m$,
$\mathfrak n$ leurs idéaux maximaux respectifs, $K=A/\mathfrak m$ et
$L=B/\mathfrak n$ leurs corps résiduels, $u:A\to B$ un homomorphisme tel que
$Bu(\mathfrak m)=\mathfrak n$, $u_0=u\otimes 1:K\to L$ l'homomorphisme
correspondant pour les corps résiduels. On considère les conditions suivantes~:
\begin{enumerate}
  \item[\rm a)] $L$ est une extension séparable de $K$ (pour $u_0$).
  \item[\rm b)] Pour tout anneau de valuation discrète complet $B'$ d'idéal
  maximal $\mathfrak n'$ et de corps résiduel $L'$,
